\chapter{I moduli in esame e i dati che producono}
\label{cap:moduli}

\begin{quote}\small
Questo capitolo risponde all'obiettivo O2. Descrive i tre sensori effettivamente
usati, il protocollo con cui si interrogano, la piattaforma di acquisizione
realizzata per la campagna sperimentale e la natura dei dati prodotti, comprese le
loro insidie --- saturazione, canali strutturalmente disattivati, coda di presenza
--- individuate in fase di validazione.
\end{quote}

\section{HLK-LD2410B}
\label{sec:ld2410b}

\subsection{Specifiche}

Il \ldb\ è un modulo radar FMCW a 24~GHz prodotto da Shenzhen Hi-Link
Electronic. Le specifiche dichiarate dal manuale ufficiale e dal datasheet sono
riassunte nella \cref{tab:ld2410b-specs}.

\begin{table*}[htbp]
\centering
\caption{Specifiche del \ldb\ (fonti: manuale ufficiale
V1.04~\cite{hilinkManualV104}, manuale V1.03~\cite{hilinkManualV103},
datasheet~\cite{ld2410bDatasheet}).}
\label{tab:ld2410b-specs}
\small
\begin{tabular}{@{}ll@{}}
\toprule
\textbf{Parametro} & \textbf{Valore} \\
\midrule
Tecnologia & FMCW, banda ISM 24~GHz \\
Portata dichiarata & 5--6~m \\
Campo visivo & $\pm 60\degree$ \\
Alimentazione & 5~V, alimentatore con capacità $>$200~mA \\
Corrente media & 80~mA \\
Interfaccia & UART 256000 baud, 8N1 \\
Bluetooth & integrato (variante B), password di default \texttt{HiLink} \\
Dimensioni & circa 35~mm $\times$ 7~mm \\
Gate di distanza & 9 (0--8) \\
Risoluzione gate & 0,75~m oppure 0,20~m \\
\bottomrule
\end{tabular}
\end{table*}

\subsection{Piedinatura e cablaggio}
\label{sec:cablaggio}

Il modulo espone cinque piedini, definiti nella Tabella~1 del
datasheet~\cite{ld2410bDatasheet}:

\begin{center}\small
\begin{tabular}{@{}clll@{}}
\toprule
\textbf{Pin} & \textbf{Nome} & \textbf{Funzione} & \textbf{Collegamento ESP32} \\
\midrule
1 & OUT & presenza digitale (alto = persona) & non collegato \\
2 & UART\_Tx & il radar trasmette & GPIO25 (RX2) \\
3 & UART\_Rx & il radar riceve & GPIO26 (TX2) \\
4 & GND & massa & GND \\
5 & VCC & alimentazione & VIN (5~V) \\
\bottomrule
\end{tabular}
\end{center}

\begin{table*}[htbp]
\centering
\caption{Corrispondenza fra i colori del cavo JST fornito con il modulo e i piedini
del datasheet. Una mappatura errata di questa corrispondenza ha comportato circa tre
settimane di lavoro perso nella convinzione che il modulo fosse guasto.}
\label{tab:colori}
\small
\begin{tabular}{@{}llll@{}}
\toprule
\textbf{Colore} & \textbf{Pin} & \textbf{Segnale} & \textbf{ESP32} \\
\midrule
Rosso & 5 & VCC & VIN (5~V) \\
Nero & 4 & GND & GND \\
Giallo & 3 & UART\_Rx & D26 (TX2) \\
Verde & 2 & UART\_Tx & D25 (RX2) \\
Blu & 1 & OUT & non collegato \\
\bottomrule
\end{tabular}
\end{table*}

\paragraph{Nota metodologica.} Vale la pena riportare l'episodio, perché è un
risultato pratico di per sé. Il 19 luglio 2026 il modulo è stato dichiarato
guasto dopo giorni di tentativi in cui non si riceveva alcun byte; il 9 agosto,
ricontrollando la Tabella~1 del datasheet invece di dedurre l'ordine dei pin dalla
fotografia del connettore, si è scoperto che i due fili della seriale erano
invertiti. Il modulo ha risposto al primo tentativo. La regola che ne deriva ---
\emph{fare sempre riferimento alla tabella di piedinatura del datasheet, non
all'ordine apparente dei fili} --- è banale ma è costata tre settimane.

Un secondo caso analogo si è presentato con lo strumento di configurazione da PC:
il pacchetto \texttt{HLK-LD2420\_TOOL} non è utilizzabile con il LD2410B (protocolli
diversi, l'errore restituito è \emph{Failed to set data transfer mode}). Per il
LD2410B va usato il \texttt{LD2410 Tool} distribuito nella cartella Drive del
produttore~\cite{hilinkDriveLD2410B}.

\subsection{Protocollo seriale}
\label{sec:protocollo-seriale}

La comunicazione avviene su UART a 256000 baud, 8N1, con due tipi di frame
distinti da intestazione e chiusura (protocollo V1.07~\cite{hilinkProtoV107}, i
dettagli completi sono in \cref{app:protocollo}):

\begin{itemize}[leftmargin=1.4em,itemsep=2pt]
  \item \textbf{frame di comando} --- intestazione \texttt{FD FC FB FA}, chiusura
        \texttt{04 03 02 01};
  \item \textbf{frame di dati} --- intestazione \texttt{F4 F3 F2 F1}, chiusura
        \texttt{F8 F7 F6 F5}.
\end{itemize}

Nel frame di dati il byte 8 contiene lo stato del bersaglio come maschera di bit
(\texttt{0x01} = in movimento, \texttt{0x02} = stazionario), seguito dalle distanze
e dalle energie del bersaglio. Se il byte 6 vale \texttt{0x01} il frame è in
\emph{engineering mode} e trasporta anche le energie per-gate nei byte 19--37, più
la lettura del sensore di luce.

I comandi utilizzati in questo lavoro sono:

\begin{center}\small
\begin{tabular}{@{}ll@{}}
\toprule
\textbf{Comando} & \textbf{Funzione} \\
\midrule
\texttt{0x00FF} & entra in modalità configurazione \\
\texttt{0x00FE} & esce dalla modalità configurazione \\
\texttt{0x0062} & attiva engineering mode \\
\texttt{0x0063} & disattiva engineering mode \\
\texttt{0x00A0} & legge la versione firmware \\
\texttt{0x0061} & legge i parametri correnti (soglie, gate, timeout) \\
\texttt{0x0060} & imposta i parametri \\
\texttt{0x00AA} & imposta la risoluzione di distanza \\
\texttt{0x000B} & auto-calibrazione del rumore di fondo (firmware V2.44) \\
\bottomrule
\end{tabular}
\end{center}

\subsection{Engineering mode e dati per-gate}
\label{sec:engineering}

In modalità normale il modulo riporta, per ogni frame, lo stato di presenza, la
distanza e l'energia del bersaglio in movimento e di quello stazionario. In
\textbf{engineering mode} aggiunge le energie di ciascuno dei nove gate,
separatamente per la componente in movimento e per quella stazionaria: diciotto
valori aggiuntivi, più la lettura del fotosensore e lo stato del pin OUT.

Questo è il punto di svolta per gli obiettivi O2, O3 e O6: senza i dati per-gate
il radar sarebbe soltanto un PIR con la distanza, mentre con essi diventa uno
strumento di misura. Tutta l'analisi del respiro e l'indice di vitalità si basano
sulle serie temporali di questi canali.

\subsection{Identità dell'esemplare in prova}
\label{sec:identita}

Per la riproducibilità dei risultati si riportano i dati letti dal modulo
effettivamente usato, ottenuti via UART il 18 agosto 2026
(\cref{tab:identita}).

\begin{table*}[htbp]
\centering
\caption{Identità e parametri di fabbrica dell'esemplare di \ldb\ in prova, letti
via UART.}
\label{tab:identita}
\small
\begin{tabular}{@{}ll@{}}
\toprule
\textbf{Dato} & \textbf{Valore letto} \\
\midrule
Versione firmware & 2.44.25070917 \\
Versione protocollo & 1 \\
MAC Bluetooth & 5E:E4:93:22:B9:3F \\
Baud confermato & 256000, 8N1 \\
Risoluzione gate & 75~cm; portata massima 675~cm (gate 0--8) \\
Timeout di presenza (\emph{no-one duration}) & 5~s \\
Soglie di movimento (gate 0--8) & 50, 50, 40, 30, 20, 15, 15, 15, 15 \\
Soglie stazionarie (gate 0--8) & 0, 0, 40, 40, 30, 30, 20, 20, 20 \\
\bottomrule
\end{tabular}
\end{table*}

Il firmware 2.44 è l'ultima versione rilasciata dal produttore ed è quella che
introduce l'\textbf{auto-calibrazione del rumore di fondo}, cioè la taratura
automatica delle soglie di movimento e stazionarie~\cite{hilinkV244}.

\paragraph{Avvertenza sulla fonte dei parametri.} Lo strumento da PC
\texttt{LD2410 Tool} riporta per le sensibilità il valore~100 su tutti i gate. È
una lettura errata: la lettura via UART con il comando \texttt{0x0061} restituisce
la scala di fabbrica decrescente riportata in \cref{tab:identita}, che è coerente
sia con il protocollo V1.07 (\S1.2.2: il bersaglio è riconosciuto solo se
l'energia supera la soglia del gate) sia con il comportamento osservato. Per i
valori numerici si è quindi fatto fede alla lettura via UART e non allo strumento
da PC.

\section{HLK-LD2420}
\label{sec:ld2420}

Il \ldc\ è il secondo radar 24~GHz in dotazione; la documentazione di
riferimento è quella distribuita dal produttore~\cite{hilinkDriveLD2420}. Le differenze rilevanti rispetto
al LD2410B sono riassunte nella \cref{tab:confronto-moduli}. Due meritano
attenzione particolare:

\begin{itemize}[leftmargin=1.4em,itemsep=2pt]
  \item \textbf{alimentazione a 3,3~V} e non a 5~V: collegarlo come il LD2410B lo
        danneggerebbe;
  \item \textbf{la piedinatura cambia con la versione firmware}. Fino alla 1.5.2 il
        piedino 3 (OT1) è l'uscita di presenza digitale e il piedino 5 (OT2) è il
        TX seriale; dalla 1.5.3 i ruoli sono invertiti, e cambia anche il baud
        (256000 fino alla 1.5.2, 115200 dalla 1.5.3). Verificare la versione prima
        di collegare non è un consiglio ma un prerequisito.
\end{itemize}

\begin{table*}[htbp]
\centering
\caption{Confronto fra i due moduli radar in dotazione.}
\label{tab:confronto-moduli}
\small
\begin{tabularx}{\linewidth}{@{}lXX@{}}
\toprule
\textbf{Caratteristica} & \textbf{LD2410B} & \textbf{LD2420} \\
\midrule
Alimentazione & 5~V & 3,3~V \\
Portata & circa 5--6~m & circa 12~m \\
Gate & 9 (0--8) & 15 (0--14) \\
Risoluzione gate & 0,75~m oppure 0,20~m & 0,70~m fisso \\
Baud & 256000 fisso & 115200 o 256000 (dipende dal firmware) \\
Bluetooth & presente & assente \\
Calibrazione & manuale (auto da V2.44) & automatica (firmware $\geq$ 1.5.4) \\
Dati per-gate & molto dettagliati & limitati \\
Piedinatura & fissa & varia con il firmware \\
Uso più adatto & respiro, dettaglio & portata lunga, copertura \\
\bottomrule
\end{tabularx}
\end{table*}

\subsection{Verifica sull'esemplare in dotazione}

L'esemplare in prova è stato collegato all'ESP32 (GPIO16/17) e verificato il
17 luglio 2026. Risultati:

\begin{itemize}[leftmargin=1.4em,itemsep=2pt]
  \item il baud funzionante è \textbf{115200}, da cui si deduce firmware
        $\geq$ 1.5.3 e quindi TX seriale sul piedino 3 (OT1). A 256000 baud si
        riceve solo rumore, con lo stesso alfabeto di byte tipico di un log letto
        al baud sbagliato;
  \item la modalità di fabbrica è \textbf{ASCII ``semplice''} e non il protocollo
        binario: il modulo trasmette righe di testo terminate da
        \texttt{CR\,LF} contenenti \texttt{ON}, \texttt{OFF} oppure
        \texttt{Range NN};
  \item in questa modalità sono disponibili solo presenza e un valore di
        ``range'': \textbf{nessuna energia per-gate}, nessuna distinzione fra
        movimento e stazionario.
\end{itemize}

\paragraph{Limite dichiarato.} L'unità di misura del campo \texttt{Range} non è
documentata nella documentazione in nostro possesso. I valori osservati
(7--37 muovendosi nella stanza) sono compatibili con decimetri, cioè 0,7--3,7~m, ma
in assenza di conferma il dato non viene usato come misura di distanza in questa
tesi. Per accedere all'engineering mode del LD2420 occorre riconfigurare il modulo
in modalità binaria; l'attività è rimasta fuori dal perimetro sperimentale
completato (\cref{sec:ld2420-risultati}).

\section{Il sensore PIR in dotazione: HC-SR501}
\label{sec:hcsr501}

Il modulo in dotazione è stato identificato per via fotografica come un
\textbf{\pir}, riconoscendo sul circuito stampato la serigrafia del circuito di
condizionamento BISS0001 e del regolatore lineare HT7133. Le specifiche da
datasheet~\cite{hcsr501} sono riassunte nella \cref{tab:pir-specs}.

\begin{table*}[htbp]
\centering
\caption{Specifiche del PIR \pir\ in dotazione (datasheet~\cite{hcsr501}).}
\label{tab:pir-specs}
\small
\begin{tabularx}{\linewidth}{@{}lX@{}}
\toprule
\textbf{Parametro} & \textbf{Valore} \\
\midrule
Base & BISS0001 + elemento LHI778 \\
Alimentazione & 4,5--20~V DC (collegato a VIN 5~V dell'ESP32) \\
Livello logico di uscita & alto 3,3~V / basso 0~V, grazie al regolatore HT7133 a
  bordo: compatibile con i GPIO dell'ESP32 senza adattamento \\
Corrente di riposo & $<$50~\textmu A \\
Portata & 3--7~m, regolabile con il trimmer di sensibilità \\
Angolo di copertura & cono $<$110\degree\ (più ampio dei $\pm 60\degree$ del radar) \\
Tempo di ritenuta & regolabile da circa 3~s a circa 5~min \\
Block time & 2,5~s dopo ogni rilascio \\
Trigger & jumper: L = singolo, H = ripetibile \\
Temperatura operativa & $-15$ / $+70$~\celsius \\
Dimensioni & 32 $\times$ 24~mm \\
\bottomrule
\end{tabularx}
\end{table*}

\paragraph{Configurazione adottata per la campagna} (fissata il 19 luglio 2026 e
mantenuta identica per tutte le acquisizioni, con fotografia della posizione dei
trimmer):

\begin{enumerate}[leftmargin=1.8em,itemsep=2pt]
  \item jumper in posizione \textbf{H} (repeat trigger), altrimenti una persona in
        movimento continuo apparirebbe intermittente per costruzione;
  \item trimmer del tempo di ritenuta al \textbf{minimo}: la ritenuta misurata è
        risultata di circa 3,4~s;
  \item trimmer di sensibilità a \textbf{metà corsa};
  \item uscita collegata a \textbf{GPIO34}. La scelta merita una nota: GPIO34 è un
        piedino di solo ingresso e privo di resistenza di pull-up interna, il che
        normalmente lo rende inadatto a un sensore a collettore aperto;
        l'\pir{} però pilota attivamente la propria uscita a 3,3~V e 0~V, quindi il
        collegamento è corretto. GPIO25 non era disponibile perché usato come RX2
        del radar.
\end{enumerate}

\section{La piattaforma di acquisizione}
\label{sec:piattaforma}

\subsection{Punto di partenza: il repository di riferimento}

Il lavoro parte dal repository del relatore~\cite{repoCallisto}, che contiene un
firmware PlatformIO per ESP32 basato sulla libreria \texttt{ncmreynolds/ld2410} e
uno script Python di acquisizione. Il firmware legge il radar via UART, gestisce
opzionalmente un display OLED e un PIR e stampa righe CSV sulla seriale; lo script
\texttt{acquire.py} le salva su file aggiungendo i metadati sperimentali passati
da riga di comando.

Lo studio del codice ha evidenziato tre limiti per gli scopi di questa tesi:

\begin{enumerate}[leftmargin=1.8em,itemsep=3pt]
  \item \textbf{campionamento a 1~Hz}. Troppo lento per misurare latenze con
        precisione e, soprattutto, insufficiente per l'analisi del respiro: con
        frequenza di campionamento 1~Hz il limite di Nyquist è 0,5~Hz, cioè
        esattamente il bordo superiore della banda respiratoria;
  \item \textbf{assenza dell'engineering mode}: viene registrata solo l'energia del
        bersaglio, non i diciotto valori per-gate;
  \item \textbf{piedinatura opposta} alla nostra (radar TX su GPIO17 e RX su
        GPIO16), quindi i due firmware non sono interscambiabili senza riscablare.
\end{enumerate}

\subsection{Il firmware realizzato}

È stato quindi scritto un logger dedicato
(\texttt{firmware/ld2410b\_logger}, \cref{app:codice}) basato sulla libreria
\texttt{MyLD2410}, con le seguenti caratteristiche:

\begin{itemize}[leftmargin=1.4em,itemsep=2pt]
  \item campionamento a \textbf{5~Hz} con periodo fisso;
  \item \textbf{engineering mode} sempre attivo, con registrazione di tutti i
        canali;
  \item lettura del PIR su GPIO34 nello stesso frame, in modo che il confronto fra i
        due sensori sia \emph{sincrono per costruzione} --- non due acquisizioni da
        allineare a posteriori;
  \item formato CSV retrocompatibile con \texttt{acquire.py}.
\end{itemize}

Il formato CSV, unico per tutta la tesi, è il seguente (le prime nove colonne
coincidono con quelle del repository di riferimento):

\begin{lstlisting}[caption={Intestazione del CSV prodotto dal logger, in
engineering mode.},label={lst:csv},basicstyle=\ttfamily\scriptsize]
timestamp_ms, radar_presence, moving_target, stationary_target,
moving_distance_cm, stationary_distance_cm, moving_energy, stationary_energy,
pir_presence,
menergy_gate0 ... menergy_gate8,   (9 colonne)
senergy_gate0 ... senergy_gate8,   (9 colonne)
light_level, out_level
\end{lstlisting}

Lo script di acquisizione aggiunge a ogni riga i metadati sperimentali:
\texttt{pc\_time\_s}, \texttt{group\_id}, \texttt{trial\_id}, \texttt{scenario},
\texttt{ground\_truth\_presence}, \texttt{ground\_truth\_state}.

\paragraph{Una correzione necessaria allo script di acquisizione.} La versione
originale scrive righe sul file solo dopo aver riconosciuto la riga di intestazione,
che l'ESP32 stampa una sola volta in \texttt{setup()}. Se il reset all'apertura
della porta seriale non scatta --- cosa che accade con una certa frequenza --- il
CSV risulta \textbf{vuoto senza alcun messaggio di errore}, e lo si scopre a
sessione finita. La nostra versione ricostruisce l'intestazione dal numero di
colonne della prima riga di dati (9 = base, 27 = engineering, 29 = engineering con
luce e OUT). È una modifica di poche righe che ha evitato la perdita di sessioni
lunghe.

\subsection{La catena di analisi}
\label{sec:catena-analisi}

L'analisi è svolta da tre script Python, scritti per questa tesi e basati sulla
sola libreria standard, tranne dove indicato:

\begin{description}[leftmargin=1.4em,style=nextline,itemsep=3pt]
  \item[\texttt{analisi/verifica\_engineering.py}] controllo di qualità di un CSV
    \emph{prima} di usarlo: cadenza reale e jitter, presenza delle colonne per-gate,
    percentuale di saturazione a 100, coerenza fra gate di picco e distanza
    riportata, rumore di fondo per-gate a stanza vuota, transizioni di presenza. Va
    eseguito a ogni sessione: è ciò che ha permesso di individuare i problemi
    descritti in \cref{sec:insidie}.
  \item[\texttt{analisi/analizza\_test.py}] metriche di confronto: accuratezza,
    falsi positivi in eventi/h, falsi negativi in percentuale, latenza di
    rilevamento e di rilascio, statistiche di distanza ed energia, aggregazione per
    scenario con media e deviazione standard sui cinque trial.
  \item[\texttt{analisi/analizza\_respiro.py}] analisi in frequenza: FFT della
    serie di energia, individuazione del picco in banda 0,1--0,5~Hz, stima degli
    atti al minuto, esportazione dello spettro in CSV. Richiede NumPy. La modalità
    \texttt{-{}-scan} prova tutti e venti i canali di energia, scarta quelli saturi
    e quelli piatti, li ordina per rapporto segnale-rumore e riporta la mediana
    delle stime concordi.
\end{description}

L'ultima funzione non era prevista inizialmente: è nata dal fatto che nel primo
tentativo di misura del respiro il canale scelto per intuito
(\texttt{stationary\_energy}) era completamente saturo, mentre il segnale era
perfettamente leggibile su \texttt{menergy\_gate2}. La ricerca automatica del
canale migliore, con criterio dichiarato, è diventata parte del metodo.

\section{Che aspetto hanno i dati: validazione e insidie}
\label{sec:insidie}

Prima di usare i dati per il confronto quantitativo è stata svolta una sessione di
validazione (673 campioni su 134~s). I risultati positivi sono:

\begin{itemize}[leftmargin=1.4em,itemsep=2pt]
  \item \textbf{cadenza di 5,00~Hz esatti}: intervallo di 200~ms su tutti i 672
        intervalli, jitter nullo. La serie temporale è uniforme, requisito
        necessario per l'analisi in frequenza;
  \item \textbf{engineering mode attivo}, con tutte e diciotto le colonne per-gate
        popolate;
  \item \textbf{mappa gate--distanza validata al 97,1\,\%}: la percentuale di
        campioni in cui il gate di massima energia coincide, entro $\pm 1$, con
        quello atteso dalla distanza riportata (distanza divisa per 0,75~m), su
        tutto il range 0--6~m. Su una seconda sessione indipendente la conferma è
        salita al 98,3\,\% su 663 campioni;
  \item \textbf{zero falsi positivi} nei 29~s di stanza vuota della sessione.
\end{itemize}

La stessa validazione ha però messo in luce tre comportamenti del modulo che
condizionano tutte le analisi successive e che vanno dichiarati esplicitamente.

\subsection{Saturazione dei canali di energia}
\label{sec:saturazione}

Con il bersaglio fermo a distanza ravvicinata, l'energia stazionaria è a fondo
scala (valore 100) nel \textbf{92\,\%} dei campioni e quella in movimento nel
\textbf{41\,\%}. Un segnale che tocca il fondo scala non porta informazione:
qualunque analisi di variazione su un canale saturo è priva di senso.

Va chiarito un punto che è stato oggetto di un'ipotesi errata iniziale: \textbf{le
soglie non risolvono la saturazione}. L'energia è un valore normalizzato 0--100 e
la soglia interviene soltanto sulla \emph{decisione} di rilevamento (protocollo
V1.07, \S1.2.2), non sul valore riportato. Nemmeno l'auto-calibrazione del firmware
V2.44 la risolve, perché anch'essa tara soglie e non la scala. L'unico rimedio è
\textbf{geometrico}: allontanare il soggetto (almeno 1,5--2~m), disallineare il
sensore o attenuare il segnale. Questo vincolo ha determinato il progetto degli
esperimenti sul respiro e sulla vitalità.

\subsection{Canali stazionari strutturalmente disattivati sotto 1,5~m}
\label{sec:senergy-zero}

Le colonne \texttt{senergy\_gate0} e \texttt{senergy\_gate1} sono risultate pari a
zero in tutti i 673 campioni. Non è un difetto dell'esemplare: anche l'esempio
riportato nel protocollo ufficiale V1.07 (\S2.3.2) mostra \texttt{00 00} per i
primi due gate stazionari, e le relative soglie di fabbrica valgono~0
(\cref{tab:identita}) --- con la regola ``energia \emph{maggiore} della soglia'' il
canale è quindi disattivato per costruzione sotto 1,5~m.

È importante non trarre la conclusione sbagliata. Il radar \textbf{riporta comunque}
bersagli fermi sotto 1,5~m: nei dati si osservano valori di
\texttt{stationary\_distance} fra 70 e 89~cm con \texttt{stationary\_energy} pari a
100. Ciò che manca è solo il \emph{dettaglio per-gate} dello stazionario nei primi
due gate. La conseguenza pratica è però rilevante: le analisi che si basano sulle
serie per-gate --- respiro e indice di vitalità --- per un soggetto a meno di 1,5~m
devono usare i \textbf{canali in movimento} oppure i gate di indice $\geq 2$.
E poiché nello scenario UPRISE la persona è sotto un banco, cioè vicina, questo è
un vincolo di progetto e non un dettaglio (\cref{sec:sotto-banco}).

\subsection{Coda di presenza}
\label{sec:coda}

Dopo l'uscita del soggetto dalla stanza il radar ha mantenuto
\texttt{radar\_presence = 1} per circa 10~s, riportando un bersaglio fermo
``fantasma'' a circa 5,2~m con energia in decadimento. Il comportamento è atteso ed
è il \emph{no-one duration} previsto dal protocollo (\S1.2.2), impostato a 5~s di
fabbrica sul nostro esemplare.

Questo ha una conseguenza metodologica precisa: la coda va misurata come
\textbf{latenza di rilascio} (\cref{sec:latenza}) e non va conteggiata come falso
positivo. Confondere le due cose porterebbe ad attribuire al radar decine di falsi
positivi inesistenti.

\subsection{Il pin OUT e il sensore di luce}

Due verifiche minori ma utili. La colonna \texttt{out\_level}, che riporta lo stato
del pin OUT letto \emph{dal frame UART}, coincide con \texttt{radar\_presence} in
1699 campioni su 1699: il collegamento fisico del filo blu è quindi inutile ai
nostri fini. È però un'informazione utile per il progetto UPRISE, perché indica che
un dispositivo interessato alla sola presenza binaria può usare il pin OUT senza
implementare la UART.

Il sensore di luce funziona: la colonna \texttt{light\_level} vale 21--29 di giorno
e 0--1 di notte, misurato su una sessione notturna di 6,5~h. Il dato è annotato
perché il comando di configurazione ausiliaria \texttt{0x01AE} non risponde sul
nostro esemplare: i due fatti sono indipendenti, il valore fotosensibile viaggia
nel frame di engineering mode e segue effettivamente l'illuminazione ambientale.

\section{Decisione sulle soglie: non ricalibrare}
\label{sec:soglie}

Un'ultima scelta metodologica va motivata, perché condiziona la validità di tutti i
confronti fra trial. Il \cref{tab:soglie} confronta le soglie di fabbrica con il
rumore di fondo misurato a stanza vuota.

\begin{table*}[htbp]
\centering
\caption{Soglie di fabbrica per-gate confrontate con il rumore di fondo misurato a
stanza vuota (29~s). Regola del protocollo V1.07 \S1.2.2: il bersaglio è
riconosciuto solo se l'energia supera la soglia.}
\label{tab:soglie}
\small
\begin{tabular}{@{}clrrrrr@{}}
\toprule
\textbf{Gate} & \textbf{Distanza} & \multicolumn{3}{c}{\textbf{Movimento}} &
\multicolumn{2}{c}{\textbf{Stazionario}} \\
\cmidrule(lr){3-5}\cmidrule(lr){6-7}
 & & soglia & rumore max & margine & soglia & rumore max \\
\midrule
0 & 0--75~cm    & 50 & 26 & 24 & 0  & 0 \\
1 & 75--150~cm  & 50 & 21 & 29 & 0  & 0 \\
2 & 150--225~cm & 40 & 9  & 31 & 40 & 6 \\
3 & 225--300~cm & 30 & 4  & 26 & 40 & 4 \\
4 & 300--375~cm & 20 & 9  & 11 & 30 & 4 \\
5 & 375--450~cm & 15 & 6  & 9  & 30 & 3 \\
6 & 450--525~cm & 15 & 11 & \textbf{4} & 20 & 5 \\
7 & 525--600~cm & 15 & 8  & 7  & 20 & 7 \\
8 & 600--675~cm & 15 & 7  & 8  & 20 & 5 \\
\bottomrule
\end{tabular}
\end{table*}

Il margine più stretto è il gate~6 in movimento (soglia 15 contro rumore massimo
11), che è anche il gate in cui si è manifestato il bersaglio fantasma della coda di
presenza. Su tutti gli altri gate il margine è ampio, e i falsi positivi osservati
sono zero. Si è quindi deciso di \textbf{non ricalibrare}: la taratura di fabbrica
regge nell'ambiente di prova, e le soglie restano il \emph{riferimento fisso} per
tutta la campagna. Cambiarle a metà campagna renderebbe i trial non confrontabili
fra loro, che è il difetto peggiore in un lavoro comparativo.

L'auto-calibrazione disponibile nel firmware V2.44 va usata dopo, come esperimento a
sé stante con confronto prima/dopo: è materiale per una sezione della tesi, non un
prerequisito della campagna.

La stabilità di questa scelta è stata verificata su tempi lunghi: nella sessione
notturna di 6,5~h la media dell'energia di rumore del gate~0 è passata da 17,4 a
17,8 fra la prima e la seconda metà della notte e quella del gate~1 da 13,3 a 13,1,
con un massimo assoluto di 34 sul gate~0 contro una soglia di 50.
